% !TeX program = uplatex
\documentclass[a4paper,12pt]{jlreq}

\usepackage{amsmath,amssymb,bm}
\usepackage{geometry}
\geometry{margin=25mm}

\begin{document}

\title{Hamiltonian Neural Networks を前提にした理論寄り研究アイデア集}
\author{}
\date{}
\maketitle

\section*{アイデア1：Hamiltonian の同定可能性（Identifiability）の理論}

\subsection*{問題設定のイメージ}

Hamiltonian Neural Networks (HNN) は軌道データから Hamiltonian
\(H(q,p)\) を学習する．しかし Hamiltonian には

\begin{itemize}
  \item 定数を足しても運動方程式は不変
  \item より一般には，正準変換（canonical transformation）で
        座標を変えても物理は同じ
\end{itemize}

という「非一意性」がある．したがって，どの意味で
「真の \(H\) が一意に決まるか」を理論的に整理することが重要である．

\subsection*{具体的なテーマ案}

位相空間が \(\mathbb{R}^2\)（1自由度系）の場合を考える．
Hamilton 方程式
\[
  \dot{z} = J \nabla H(z), \qquad
  z = (q,p),\quad
  J = \begin{pmatrix} 0 & 1 \\ -1 & 0 \end{pmatrix}
\]
を満たすベクトル場が観測されるとき，2つの Hamiltonian
\(H_1, H_2\) が同じベクトル場を与える条件を調べる：
\[
  J \nabla H_1(z) = J \nabla H_2(z)
  \quad\Rightarrow\quad
  \nabla(H_1 - H_2)(z) = 0
  \quad\Rightarrow\quad
  H_1 - H_2 = \text{const}.
\]

\begin{itemize}
  \item 有限時間区間内の軌道データ（ノイズあり）の下で，
        \(H(q,p)\) がどのノルムでどこまで一意に決まるかを評価する．
  \item 単純なクラス（例：分離形
        \(H(q,p)= \dfrac{p^2}{2m} + V(q)\) で \(V\) が滑らか）に限定して，
        軌道情報から \(V(q)\) を一意に同定できる条件を定理として与える．
\end{itemize}

\subsection*{HNN 損失との接続}

HNN の典型的な損失は
\[
  \mathcal{L}(\theta)
  =
  \bigl\|
    \dot{z}_{\text{data}} - J \nabla H_\theta(z)
  \bigr\|^2
\]
となる．このとき，\(\mathcal{L}(\theta)\) が十分小さいとき，
「真の Hamiltonian」\(H^\star\) と \(H_\theta\) の差
\(\| H_\theta - H^\star\|\) をどの程度まで抑えられるかを評価する
簡易版の誤差評価を与えることが目標となる．

\bigskip

\section*{アイデア2：学習誤差付き HNN の長時間ダイナミクス解析}

\subsection*{設定}

真の Hamiltonian を \(H^\star\)，学習されたものを
\[
  H_\theta = H^\star + \varepsilon R
\]
と書き，\(\varepsilon \ll 1\) を学習誤差の大きさとみなす．
連続時間系は
\[
  \dot{z} = J \nabla H^\star(z),
  \qquad
  \dot{z}_\theta = J \nabla H_\theta(z_\theta)
\]
となる．さらに数値積分として，例えばシンプレクティック積分
（Verlet 法など）を仮定する．

\subsection*{解析のゴール（1自由度系に絞る）}

\begin{enumerate}
  \item 単振り子や調和振動子などの 1 自由度系について，
        \(\varepsilon\) を小パラメータとみなした摂動解析により
        \[
          z_\theta(t) - z^\star(t)
        \]
        の成長を評価する．
  \item 数値時間刻み \(h\) を導入したとき，
        エネルギー誤差
        \[
          \Delta H(t)
          =
          H^\star\bigl(z_{\text{num}}(t)\bigr)
          - H^\star\bigl(z^\star(t)\bigr)
        \]
        の時間スケーリング（例えば \(\mathcal{O}(\varepsilon)\),
        \(\mathcal{O}(h^2)\) など）を評価する．
  \item ベンチマークとして，ベクトル場
        \(\dot{z} = f_\phi(z)\) を直接ニューラルネットで近似する
        「非構造的」モデルと比較し，
        長時間安定性（エネルギーのドリフト）に関する
        理論・数値両方の評価を行う．
\end{enumerate}

\subsection*{理論的背景}

\begin{itemize}
  \item 連続 Hamilton 系に対するシンプレクティック積分法の
        長時間保存性（バックワードエラー解析）．
  \item Hamiltonian の摂動論（簡易版 KAM 的な議論の 1 自由度版）．
\end{itemize}

完全な KAM 理論ではなく，1 自由度に限定して線形化＋摂動解析を行うことで，
学部レベルでも手が届く設定にするのが現実的である．

\bigskip

\section*{アイデア3：HNN の表現力を単純なクラスで完全解析}

\subsection*{対象とするクラス}

1 自由度で分離形の Hamiltonian を考える：
\[
  H(q,p)
  =
  \frac{p^2}{2m} + V(q),
  \qquad
  V \in C^k([a,b]) .
\]
このクラスに対して，HNN（もしくは類似構造のネットワーク）が
どの程度の規模（層の深さ・幅）で \(\varepsilon\)-近似できるかを
明示的に評価することを目指す．

\subsection*{基本方針}

\begin{enumerate}
  \item ReLU ネットワークや Tanh ネットワークによる
        一変数関数 \(V(q)\) の近似定理を用いて，
        \(\|V - V_\theta\|_\infty \le \varepsilon\) を達成するための
        ネットワークサイズの上界を与える．
  \item その誤差を Hamiltonian 全体に拡張し，
        ベクトル場
        \[
          f(z) = J \nabla H(z)
        \]
        の近似誤差 \(\| f - f_\theta \|\) を見積もる．
  \item 誤差 \(\varepsilon\) を達成するための
        パラメータ数，層数，幅のスケーリングを
        明示的なオーダー（例：\(\mathcal{O}(\varepsilon^{-\alpha})\)）
        で与えることを目標とする．
\end{enumerate}

\subsection*{既存研究との差別化}

既存のユニバーサル近似定理は，より一般の Hamiltonian 写像や
高次元系を対象とした抽象的な結果が多く，
ネットワーク規模の評価も漸近的・存在的なものが中心である．

ここでは対象クラスを
\begin{itemize}
  \item 1 自由度
  \item 分離形 Hamiltonian
\end{itemize}
に強く制限する代わりに，より具体的で鋭いサイズ評価を与えることが
新規性となる．

\bigskip

\section*{アイデア4：正準変換に対する不変性と HNN}

\subsection*{問題意識}

Hamiltonian 力学では，正準変換
\[
  (q,p) \mapsto (Q,P)
\]
の下でシンプレクティック形式が保存される．
HNN で学習された \(\hat{H}(q,p)\) が，正準変換後の座標
\((Q,P)\) でどのように表現されるべきか，
またそれがどのような意味で「同じ物理系」を表すかを
理論的に整理する．

\subsection*{1自由度での具体的検討}

\begin{itemize}
  \item 調和振動子など単純な系に対して，
        \((q,p)\) からアクション--アングル変数 \((J,\theta)\) への
        正準変換を具体的に書き下ろす．
  \item \((q,p)\) 座標で学習した Hamiltonian \(\hat{H}(q,p)\) と，
        \((J,\theta)\) 座標で学習した \(\tilde{H}(J,\theta)\) の関係を
        理論的・数値的に解析し，
        どの条件のもとで
        \[
          \hat{H}(q,p) \simeq H^\star(q,p) + \text{const.},\quad
          \tilde{H}(J,\theta) \simeq K(J) + \text{const.}
        \]
        といえるのかを議論する．
\end{itemize}

\subsection*{論文構成イメージ}

\begin{enumerate}
  \item \textbf{導入・背景}：  
        Hamiltonian 力学と正準変換の基礎，
        HNN・Symplectic NN の概要と既存研究の整理．
  \item \textbf{理論部分}：  
        1 自由度 Hamilton 系における正準変換のもとでの
        Hamiltonian の変換則と，
        HNN による学習結果の「同値性」の定義・命題の提示．
  \item \textbf{簡単な数値例}：  
        調和振動子などに対し，
        2 つの座標系で HNN を学習し，
        学習された Hamiltonian の関係を検証する．
\end{enumerate}

このアイデアはシンプレクティック幾何と HNN の橋渡しをする
「概念整理＋簡単な定理＋数値例」という構成であり，
理論寄りの卒論・修論としてまとめやすいテーマとなる．









%==============================
% 教科書PDFとHNN論文の対応
%==============================

\section*{付録：教科書PDFと HNN 論文の対応関係}

本節では，配布されている教科書 PDF（第13章）と
Greydanus らの Hamiltonian Neural Networks (HNN) 論文との
対応関係を整理する．

\subsection*{13.1～13.4：仮想仕事～変分原理と論文の対応}

教科書 PDF の前半（13.1～13.4）には
\begin{itemize}
  \item 13.1 仮想仕事の原理
  \item 13.2 ダランベールの原理
  \item 13.3 ラグランジュ方程式の導出 その1
  \item 13.4 変分原理（最小時間・フェルマー原理など）
\end{itemize}
が含まれている．

HNN 論文では，これらの内容は直接は登場しないが，
\emph{ニュートン力学 \(\to\) ラグランジュ形式 \(\to\) ハミルトン形式}
という流れを理解するための前提知識となる．

論文の 2. Theory の冒頭では，
古典力学の複数の表現形式をさらっと触れつつ，
「本論文ではハミルトン形式だけを用いる」という立場が取られているため，
その背景を理解するために 13.1～13.4 を読んでおくとよい．

\subsection*{13.5～13.7：ハミルトンの原理・一般化座標と論文の対応}

PDF の次の部分（13.5～13.7）には
\begin{itemize}
  \item 13.5 ハミルトンの原理
  \item 13.6 ラグランジュの運動方程式の導出 その2
  \item 13.7 一般化座標と一般化力
\end{itemize}
が含まれている．

HNN 論文 2. Theory の
\emph{Hamiltonian Mechanics} 節では，
状態を一般化座標 \(q\) とその共役運動量 \(p\) の組
\((q,p)\) で表すとして議論を始める．これは教科書第13章における
\begin{itemize}
  \item 一般化座標 \(q_k\)
  \item 一般化運動量 \(p_k\)
\end{itemize}
の導入に対応している．

したがって，
\[
  (q,p) \text{ を状態ベクトルとして扱うとはどういうことか}
\]
をきちんと理解したい場合には，13.7（およびその前後）を読んでおくと
論文の数式が自然に受け入れられる．

\subsection*{13.8～13.9：ラグランジュ方程式と具体例 ↔ 論文のタスク系}

PDF の 13.8～13.9 では
\begin{itemize}
  \item 13.8 一般化座標で表したラグランジュの運動方程式
  \item 13.9 ラグランジュ方程式の応用（単振り子，円周上質点，二重振り子など）
\end{itemize}
が扱われている．

HNN 論文 3. Learning a Hamiltonian from Data では，
以下のような物理系をタスクとして扱っている：
\begin{itemize}
  \item 質点バネ系（ideal mass-spring）
  \item 単振り子（ideal pendulum）
  \item 2体問題（two-body problem） など
\end{itemize}
これらの系の運動方程式を，自分の手で
ラグランジュ形式から導けるようになりたいなら，
教科書の 13.8～13.9 が直接対応する．

また HNN 論文では
\[
  H = K + U
\]
という形で Hamiltonian を
運動エネルギー \(K\) とポテンシャルエネルギー \(U\) の和として書いているが，
これは教科書 13.8～13.9 で導入される
「エネルギー分解 \(K, U\)」の考えそのものである．

\subsection*{13.10：ハミルトンの正準方程式 ↔ 論文の式 (2)}

教科書 13.10「ハミルトンの正準方程式」では，
ラグランジアン \(L(q,\dot q,t)\) から正準運動量 \(p_k\) を定義し，
レジェンドル変換により Hamiltonian \(H(q,p,t)\) を導入したうえで，
\[
  \dot q_k = \frac{\partial H}{\partial p_k},
  \qquad
  \dot p_k = -\frac{\partial H}{\partial q_k}
\]
というハミルトンの正準方程式を導いている．

HNN 論文 2. Theory の中核となる式 (2) は，まさにこれに対応している：
\[
  \dot{q}, \dot{p} = J \nabla H(q,p),
\]
ただし \(J\) は
\[
  J =
  \begin{pmatrix}
    0 & I \\
    -I & 0
  \end{pmatrix}
\]
のような交代行列である．

HNN はニューラルネットワークによりスカラー関数 \(H_\theta(q,p)\) を出力し，
その勾配 \(\nabla H_\theta\) を用いて
\[
  \dot z_\theta = J \nabla H_\theta(z)
\]
を計算し，これをデータ \(\dot z_{\text{data}}\) と比較する損失
\[
  \mathcal{L}(\theta)
  =
  \bigl\|
    \dot{z}_{\text{data}} - J \nabla H_\theta(z)
  \bigr\|^2
\]
で学習する．

したがって，
\begin{quote}
  HNN 論文の「Hamiltonian Mechanics」節と HNN の定義部分は，
  教科書 13.10 の内容をそのまま利用している
\end{quote}
と見なしてよい．

\subsection*{13.11：ポアンソンの括弧 ↔ Pixel Pendulum の Poisson bracket}

教科書 13.11 ではポアンソン括弧
\[
  [F,G]
  =
  \sum_k \left(
    \frac{\partial F}{\partial q_k}\frac{\partial G}{\partial p_k}
    -
    \frac{\partial F}{\partial p_k}\frac{\partial G}{\partial q_k}
  \right)
\]
が定義され，その性質が紹介される．

HNN 論文の Pixel Pendulum の実験部分では，
オートエンコーダの潜在ベクトル \(z\) の前半を「座標」，
後半を「その共役量（あるいは時間微分）」とみなし，
ポアンソン括弧が正準座標に近い性質
\[
  [q_i,p_j] \approx \delta_{ij},
  \quad
  [q_i,q_j] \approx 0,
  \quad
  [p_i,p_j] \approx 0
\]
を満たすように補助損失を設計している．

このときに用いられる \emph{Poisson bracket relations} が，
まさに教科書 13.11 の内容に対応している．

\subsection*{どの順番で読めば HNN が理解しやすいか}

HNN 論文を理解するうえで，特に直接役に立つ教科書の節は
\begin{enumerate}
  \item 13.7 一般化座標と一般化力
  \item 13.10 ハミルトンの正準方程式
  \item 13.9（単振り子や2体問題などの具体例）
  \item 13.11 ポアンソンの括弧
\end{enumerate}
である．

したがって，例えば次の順で読むと，
HNN 論文の数式部分がかなりスムーズに追えるようになる：
\[
  13.7 \;\to\; 13.10 \;\to\; 13.9 \;\to\; 13.11.
\]

この付録を参照しながら教科書と論文を行き来することで，
\begin{itemize}
  \item 解析力学としての理論的理解
  \item HNN の数式の意味づけ
\end{itemize}
の両方を強化できる．


\end{document}

